\section{Scaling Law 与模型架构}

\subsection{预训练 Scaling 的新瓶颈}

\begin{warningbox}{高质量数据接近常数}
高质量训练数据的增长速度远慢于算力增长：
\begin{itemize}
    \item 高质量文本数据量可以近似看作常数
    \item 多模态数据无法有效提升文本的"智商"
    \item 这使得 Token Efficiency 成为比 Compute Efficiency 更重要的优化方向
\end{itemize}
\end{warningbox}

\subsection{Muon 优化器：突破 Adam 的 10 年统治}

Adam 优化器统治了深度学习训练领域近 10 年。Muon 优化器通过考虑参数之间的矩阵结构依赖关系，实现了约 2 倍的 Token Efficiency 提升：

\begin{itemize}
    \item Adam：独立考虑每个参数元素，忽略参数间的结构关系
    \item Muon：考虑矩阵参数间的 dependency，通过正交化方式优化
    \item 效果：30T 高质量 Token 的学习效果等价于 Adam 学习 60T Token
\end{itemize}

\subsection{RL Scaling 的发现}

\begin{knowledgebox}{强化学习中的 Curriculum 设计}
在 Agent 训练中，不能一上来就让模型解决最难的任务（如证明未解数学猜想），否则 Sample Efficiency 极低。更好的方式是：
\begin{itemize}
    \item 使用中等难度的任务进行训练
    \item 逐步提升难度（爬山过程）
    \item 搭配好的 Sampling 策略
\end{itemize}
这本质上是一个 Curriculum Learning 的过程。
\end{knowledgebox}

\subsection{本章小结}

预训练 Scaling 的新瓶颈是高质量数据的有限性。Muon 优化器通过矩阵正交化突破 Adam 的 10 年统治，实现 2 倍 Token Efficiency 提升。RL 训练需要 Curriculum 设计来保证 Sample Efficiency。


